%! Problems with Sampling — Unit 3, Lesson 3.3 — Guided Notes & Practice
% THE in-class document, in the MAIN IDEAS / NOTES shape (see LESSON_SHAPE.md §1), at the
% COMPACT budget: the vocabulary box, then ONE guidednotes table of FOUR rows (three instruction
% rows + Guided Practice), 14 problems.
%   I DO   : the teacher fills the definition lines and works the FIRST problem of each grid
%            (1, 4, 7).
%   WE DO  : the class works the rest; the last row, Guided Practice, is worked together.
%   YOU DO : nothing on this page — ../homework/ IS the individual practice.
% THE AP HABIT, every problem: name the bias, say who is over- or under-represented and why,
% and give the DIRECTION (too high / too low) in context.
% THE TRAP is problem 3 (comment cards: voluntary response bias pushes the estimate DOWN —
% the direction comes from who cares enough to answer, not from the name of the bias).
% THE CRUX is problem 9, in the last instruction row: 48,000 forum replies "must be close to
% the truth" — a larger sample reduces variability, not bias (EK DAT-2.E.2, DAT-2.E.6).
%
% THE CONTEXTS (instruction, guided practice, and homework each use a different one)
%  Row 1 — one-line surveys: a news-site traffic poll, an AP Calculus class, comment cards.
%  Row 2 — a landline phone book, a mailed work-hours survey, after-school club rosters.
%  Row 3 — a turf-field question, self-reported study time, a video-game forum survey (crux)
%          against 600 players emailed at random.
%  Row 4 (GP) — a university's survey of students' paid work hours: dorm list, advisor
%          meetings, a face-to-face question, and "email all 30,000".
%
% PAGE LOCKSTEP with notes_key/: \termblank -> \termans, \blank -> \ans, 4th arg of \pcell.
% TABLE MECHANICS: inside a cell `\\` ENDS THE ROW — use \par. Long rows are split into
% sub-rows (a bare `\\` then `& ...`) so the table can break between them.
% NO SKETCH QUESTIONS: every display is given and read.
\documentclass[12pt]{article}
\usepackage{apstats-article}
\usepackage{apstats-boxes}

\begin{document}

\pageheader{Unit 3, Lesson 3.3}{Guided Notes \& Practice}
% NAMESTRIP: no \namedateperiod here — the name row lives on the cover only.

\vspace{2pt}

\boxguard[16]
\begin{vocabbox}
\small
Fill in each term \textbf{as we name it} in the notes below.
\par
\termblank{Voluntary response bias}
\par
\termblank{Undercoverage bias}
\par
\termblank{Nonresponse bias}
\par
\termblank{Response bias}
\par
\end{vocabbox}

\small
\begin{guidednotes}
% ── ROW 1: who chose to be in the sample? — THE TRAP (problem 3) ─────────────
\mainidea[Who chose?]{Volunteers} &
In a \textbf{voluntary response sample}, people \blank{1.7cm} themselves, and those with
\blank{1.4cm} opinions answer most, so the result leans their way: \textbf{voluntary response
bias}. A \textbf{convenience sample} takes whoever is \blank{1.7cm} to reach. Neither uses
\blank{1.4cm} to pick who is in the sample.

\vspace{3pt}
\textbf{A complete answer, every time:}\par
\stepnum{1}\ Name the bias.\par
\stepnum{2}\ Say who is over- or under-represented, and why, in context.\par
\stepnum{3}\ Give the \blank{1.9cm}: is the estimate too high or too low?
\notesprompt{Name the bias. Too high or too low? Why?}
\begin{probgrid}{|Y|Y|Y|}
\pcell{1}{A news site asks, ``Is traffic in our town a serious problem?'' 79\% of 2{,}100 readers who vote say yes.}{1.9cm}{} &
\pcell{2}{To estimate nightly homework time at her school, Ana asks the 30 students in her AP Calculus class.}{1.9cm}{} &
\pcell{3}{A restaurant leaves comment cards on its tables. 40 of 600 diners return one; 55\% say they were satisfied.}{1.9cm}{} \\ \hline
\end{probgrid}
\\ \hline
% ── ROW 2: who was left out, and who never answered? ─────────────────────────
\mainidea[Who is missing?]{Coverage and response} &
\textbf{Undercoverage}: part of the population has a \blank{1.7cm} chance --- or no chance
--- of being chosen, often because the \blank{1.0cm} we sample from leaves them out.
\textbf{Nonresponse}: individuals \blank{1.6cm} for the sample can't be reached or
\blank{1.4cm}, and they may differ from those who answer. Either one can spoil a
\blank{1.6cm} sample.

\vspace{3pt}
\textbf{The test:} Could they ever have been picked? \emph{No} $\to$ undercoverage.
\emph{Picked, but no answer} $\to$ nonresponse.
\notesprompt{Name the bias. Too high or too low? Why?}
\begin{probgrid}{|Y|Y|Y|}
\pcell{4}{To estimate residents' hours online per day, a town calls 500 numbers picked at random from the landline phone book.}{1.9cm}{} &
\pcell{5}{A city mails a long survey on weekly hours worked to 1{,}000 random addresses; 310 come back.}{1.9cm}{} &
\pcell{6}{To estimate the percent of students with a part-time job, a principal picks 80 at random from the after-school club rosters.}{1.9cm}{} \\ \hline
\end{probgrid}
\\ \hline
% ── ROW 3: how the question was asked — and THE CRUX (problem 9) ─────────────
\mainidea[How was it asked?]{Response bias} &
\textbf{Response bias} comes from \blank{1.0cm} the data are gathered, not from who is chosen:
a \blank{1.6cm} or confusing question, or an answer people won't give
\blank{1.7cm} about themselves (self-reported data).
\notesprompt{Name the bias. Too high or too low? Why?}
\begin{probgrid}{|Y|Y|}
\pcell{7}{Students are asked, ``Would you support a new turf field so athletes can play safely in any weather?'' 74\% say yes.}{1.6cm}{} &
\pcell{8}{A teacher asks each student, out loud, ``How many hours did you study for the test?''}{1.6cm}{} \\ \hline
\end{probgrid}
\\
& A game company posts a survey about its new update on its forum: \textbf{48{,}000} players
reply and \textbf{81\%} dislike the update. It also emails \textbf{600} players picked at random
from all 3 million accounts: \textbf{47\%} dislike it.
\begin{probgrid}{|Y|Y|}
\pcell{9}{``48{,}000 replies is so many that 81\% must be close to the truth for all players.'' Agree? Why?}{1.6cm}{} &
\pcell{10}{Which estimate would you report for all 3 million players? Why?}{1.6cm}{} \\ \hline
\end{probgrid}
\par\vspace{4pt}
\fcolorbox{navy}{sky}{\parbox{\dimexpr\linewidth-2\fboxsep-2\fboxrule\relax}{\centering
\textbf{A larger sample reduces \blank{2.4cm}, not \blank{1.3cm}.}}}
\\ \hline
% ── ROW 4: guided practice — WE DO, together, fresh context ──────────────────
\mainidea[Guided practice]{The Work-Hours Survey} &
A university wants the mean number of hours per week its students work at paid jobs.
\par\vspace{3pt}
\begin{minipage}{\linewidth}\centering\footnotesize
\setlength{\tabcolsep}{4pt}%
\renewcommand{\arraystretch}{1.2}
\begin{tabular}{@{} c p{0.82\linewidth} @{}}
\rowcolor{sky}
\textbf{Step} & \textbf{The plan} \\
\hline
1 & Pick 500 students at random from the list of students who live in campus dorms. \\
2 & Ask each to meet with their academic advisor; only 140 show up. \\
3 & In the meeting, the advisor asks, ``How many hours a week do you work at a paid job?'' \\
\end{tabular}
\end{minipage}
\notesprompt{We work these together.}
\begin{probgrid}{|Y|Y|}
\pcell{11}{Step 1: who can never be picked? Name the bias. Too high or too low?}{1.6cm}{} &
\pcell{12}{Step 2: only 140 of 500 come. Name the bias. Which way? Why?}{1.6cm}{} \\ \hline
\pcell{13}{Step 3: why might the answers be off? Name the bias and its direction.}{1.6cm}{} &
\pcell{14}{``Next year we'll email all 30{,}000 students and get 6{,}000 replies. That fixes the bias.'' Respond.}{1.6cm}{} \\ \hline
\end{probgrid}
\\ \hline
\end{guidednotes}

% ── THE NOTES END AT GUIDED PRACTICE ─────────────────────────────────────────
% There is deliberately no "you do" here: ../homework/ is the individual practice,
% started in class in the last 8 minutes while the teacher circulates.

\end{document}
